\documentclass{article}
\usepackage[margin=1in]{geometry}
\usepackage{graphicx}
\usepackage{float}
\usepackage{microtype}
\usepackage{textcomp, gensymb}
\usepackage{enumitem}
\usepackage{amsmath}
\usepackage{parskip}
\begin{document}
\setcounter{section}{4}
\section{Chapter 5: Magnetostatics}
\setcounter{subsection}{-1}
\subsection{Overview}
In electrostatics, our primary concern was the relationship between electric field, \(\vec{E}\), and electric flux density, \(\vec{D}\), whose relationship is defined as
\[
  \vec{D} = \varepsilon \vec{E},
\]
where \(\varepsilon\) refers to the electric permittivity of the material. 

When discussing \textit{magnetic} fields, the subject of this chapter, we will consider an analogous relationship described by 
\[
  \vec{B} = \mu \vec{H},
\]
where \(\vec{B}\) is the magnetic flux density, \(\vec{H}\) is the magnetic field intensity, and \(\mu\) is the medium's \textbf{magnetic permeability}. Unless otherwise specified, we can assume that the magnetic permeability of most materials is \(\mu_0 = 4 \pi \times 10^{-7} \ \mathrm{\frac{H}{m}}\).

A key difference between electric and magnetic fields is that whereas electric fields are induced by static charges, magnetic fields are generated by \textit{moving} charges; or current. 

\subsection{Magnetic Forces and Torques}
Just as magentic fields are \textit{generated} by moving charges, only moving charges experience a force when in a magnetic field. The resultant force may be described as 
\[
	F_{\mathrm{m}} = q \vec{u} \times \vec{B},
\]
	where \(\vec{u}\) is the affected charge's velocity. Note that, using the definition of the cross product, the induced force will be at its maximum when \(\vec{u}\) and \(\vec{B}\) are perpendicular, and \(0\) when they are parallel.

In the case of moving charges, we must consider the possibility of both the influence of magnetic and electric fields on the point charge. The resultant total force on the point charge is simply the sum of the two forces, called the Lorentz Force equation,
\[
	\vec{F} = \vec{F}_e + \vec{F}_m = q \vec{E} + q \vec{u} \times \vec{B} = q(\vec{E}+\vec{u} \times \vec{B}) \ \mathrm{N}.
\]

Since moving point charges experience a force in a magnetic field, we may generalize this effect to that current carrying conductors, which are essentially collections of moving charges, and also experience magnetic forces under the same circumstances. Mathematically, we may integrate the total influence of the charges in a conductor; assuming a closed circuit (the only kind in which constant current can flow), 
\[
	\vec{F}_{\mathrm{m}} = I \oint_C d \vec{l} \times \vec{B} \ \mathrm{N}.
\]
 
As an example, given a semicircular conductor carrying a current \(I\) exposed to a magnetic field \(\vec{B} = \hat{\alpha}_y B_0\), find:
\begin{enumerate}[label=\alph*)]
	\item the magnetic force \(\vec{F}_1\) on the straight section of the wire.

		For the straight section (bottom of the semicircle) we may define \(d \vec{l} = \hat{\alpha}_x dx\), since the entire section is in the \(x\) direction. Then, 
		\begin{align*}
			&\vec{F}1 = I \int d \vec{l} \times \vec{B} = I \int \hat{\alpha}_x dx \times \hat{\alpha}_y B_0 \\
			&\vec{F}_1 = IB_0\hat{\alpha}_z \int_{0}^{2r}dx = \hat{\alpha}_z 2IrB_0 \ \mathrm{N}.
		\end{align*}
	\item The magnetic force \(\vec{F}_2\) on the curved section.
		Based on the cross product \(d \vec{l} \times \vec{B}\), we know that the force will be in the \(-\hat{\alpha}_z\) direction (right-hand rule). Finding \(d \vec{l}\), using the cylindrical coordinate system, \(d \vec{l} = r d\phi \hat{\alpha}_\phi\). Then, 
		\begin{align*}
			&\vec{F}_2 = I \int d \vec{l} \times \vec{B} = I \int r d \phi \hat{\alpha}_\phi \times \hat{\alpha}_y B_0\\
			&\vec{F}_2 = I r B_0 \int (-\sin{(\phi)} \hat{\alpha}_x + \cos{(\phi)} \hat{\alpha}_y) \times \hat{\alpha}_y \ d \phi \\
			&\vec{F}_2 = I r B_0 (-\hat{\alpha}_z) \int_{0}^{\pi} \sin{(\phi)} d \phi = -\hat{\alpha}_z 2 I r B_0 \ \mathrm{N}.
		\end{align*}
\end{enumerate}
Therefore, \(\vec{F}_1 = -\vec{F}_2\), and there is no net force. 

Current carrying loops under the influence of magnetic fields can also \textit{rotate} due to magnetic torque. For example, given a magnetic field in the \(x-y\) coordinate system \(\vec{B} = \hat{\alpha}_y B_y+ \hat{\alpha}_z B_z\). Given a current carrying loop (square) in the \(x-y\) plane, what current direction will produce torque under the given magnetic field? Testing all the available options, 
\begin{align*}
	& I_1 = \hat{\alpha}_y I\\
	& I_2 = -\hat{\alpha}_x I\\
	& I_3 = -\hat{\alpha}_y I\\
	& I_4 = \hat{\alpha}_x I.
\end{align*}
Then, using \(\vec{F}_m = \int_{L} (\vec{I} \times \vec{B}) \ dl\),
\begin{align*}
	& \vec{F}_1 = (\hat{\alpha}_y I \times \hat{\alpha}_z B_z) l_1 = \hat{\alpha}_x I B_z l_1\\
	& \vec{F}_2 = (-\hat{\alpha}_x I \times \hat{\alpha}_z B_z) l_2 = \hat{\alpha}_y I B_z l_1\\
	& \vec{F}_3 = (-\hat{\alpha}_y I \times \hat{\alpha}_z B_z) l_1 = -\hat{\alpha}_x I B_z l_1\\
	& \vec{F}_4 = (\hat{\alpha}_x I \times \hat{\alpha}_z B_z) l_2 = -\hat{\alpha}_y I B_z l_2\\
\end{align*}
Therefore, for the normal (\(B_z\)) component of the field, the loop is pulled in all directions equally, and no torque results. Finding the effect of the \(B_y\) components, 
\begin{align*}
	& \vec{F}_1 = (\hat{\alpha}_y I \times \hat{\alpha}_y B_y) l_1 = 0 \\
	& \vec{F}_2 = (-\hat{\alpha}_x I \times \hat{\alpha}_y B_y) l_2 = -\hat{\alpha}_z I B_y l_1\\
	& \vec{F}_3 = (-\hat{\alpha}_y I \times \hat{\alpha}_y B_y) l_1 = 0 \\
	& \vec{F}_4 = (\hat{\alpha}_x I \times \hat{\alpha}_y B_y) l_2 =  \hat{\alpha}_x I B_y l_2.
\end{align*}
Therefore, net force is still zero, but now, since \(\vec{F_1}\) and \(\vec{F_2}\) are equal but in opposite directions, there will be a torque that causes rotation in \textit{what direction}? 

More generally, torque can be found as \(\vec{m} = IA \hat{\alpha}_n\), where \(\vec{m}\) is the vector magnetic moment, \(A\) is the loop area, and \(\hat{\alpha}_n\) is the unit vector normal to the loop (right-hand rule applied to current direction). Then, 
\[
  \vec{T} = \vec{m} \times \vec{B},
\]
where \(\vec{T}\) is the direction of torque. Alternately, 
\begin{align*}
	& \vec{T} = \hat{\alpha}_T m B \sin{(\theta)} \\
	&\vec{T} = \hat{\alpha}_T I A B \sin{(\theta)}.
\end{align*}
\subsection{The Biot-Savart Law}
The Biot-Savart law defines the static magnetic field produced by a steady current. 

To formulate the Biot-Savart law, we take a short section \(dl\) of a current carrying wire, and define the resultant magentic field at a point lying at a straight-line distance \(R\) from the current carrying wire, 
\[
	d \vec{H} = \frac{I}{4\pi}\frac{d \vec{l} \times \hat{\alpha}_r}{R^2}.
\]
Integrating over the entire current carrying wire, 
\[
  \vec{H} = \int d \vec{H} = \frac{I}{4\pi} \int \frac{d \vec{l}\times \hat{\alpha}_r}{R^2},
\]
where \(\hat{\alpha}_r = \frac{\vec{R}}{|\vec{R}|}\).

Alternately, 
\[
\vec{H} = \frac{I}{4\pi} \int \frac{d \vec{l} \times \vec{R}}{R^3} \ \mathrm{\frac{A}{m}}.
\]
The direction of the resultant field will be rotating in the \(\phi\) direction around the wire. Note that \(d \vec{l}\) is in the direction of the current \(I\).

Solving problems using Biot-Savart:

For example, we have a linear conductor of length \(l\) carrying a current of \(I\) along the \(z\) axis. Find the magnetic flux density \(\vec{B}\) at a point \(P\) located at a distance of \(r\) in the \(x-y\) plane.

To make the problem simpler, we may center the conductor at the origin, meaning one end is at \(-\frac{l}{2}\) and the other is at \(\frac{l}{2}\). We know that the magnetic field will be in the \(\hat{\alpha}_\phi\) direction, and we assume a magnetic permeability of \(\mu_0\). Therefore, 
\[
  \vec{H} = \frac{I}{4\pi} \int \frac{d \vec{l} \times \vec{R}}{R^3}.
\]

First, we need to define a few parameters. In this case, \(d \vec{l} = \hat{\alpha}_z dz\), since the conductor exists along the \(z\) axis. Finding the difference vector between any point on along the conductor and \(P\), we can express \(\vec{R}\) as \(\hat{\alpha}_r \cdot r - \hat{\alpha}_z \cdot z\). Then, \(d \vec{l} \times \vec{R} = (\hat{\alpha}_z dz) \times (\hat{\alpha}_r \cdot r - \hat{\alpha}_z \cdot z) = \hat{\alpha}_\phi r dz\).

Then, 
\[
	\vec{H} = \frac{I}{4\pi}\int_{-\frac{l}{2}}^{\frac{l}{2}} \frac{\hat{\alpha}_\phi r dz}{(z^2 + r^2)^{\frac{3}{2}}} = \hat{\alpha}_\phi \frac{I}{4\pi}\cdot r \int_{-\frac{l}{2}}^{\frac{l}{2}} \frac{dz}{(z^2 + r^2)^{\frac{3}{2}}}.
\]
Using a standard formula for integration, this reduces to
\[
	\vec{H} = \hat{\alpha}_\phi \frac{I\cdot r}{4\pi} \cdot \frac{I}{r^2}(\frac{z}{\sqrt{r^2+z^2}} \ \Big|_{-\frac{l}{2}}^{\frac{l}{2}}) = \hat{\alpha}_\phi \frac{I\cdot l}{2\pi r \sqrt{4r^2 + l^2}}.
\]

Finding \(\vec{B}\), 
\[
	\vec{B} = \mu_0 \vec{H} = \hat{\alpha}_\phi\frac{\mu_0 I \cdot l}{2\pi r \sqrt{4r^2 + l^2}} \ \mathrm{T}.
\]

We may simplify this equation for an identical conductor to the above conductor, except that it extends to infinity. So, as \(l \to \infty\),
\[
  \vec{B} = \hat{\alpha}_\phi \frac{\mu_0 I}{2\pi r}.
\]

Later, we will show that the same problem can be solved more easily using Ampere's Law, the magnetic field analogue to Gauss's Law.

As another example, a circular loop of radius \(a\) carries a current of \(I\). Find the magnetic flux density \(\vec{B}\) at a point on the axis of the loop. 

Use the right-hand rule to determine that, if the current \(I\) flows counter-clockwise, the field will be in the \(\hat{\alpha}_z\) direction (fingers are current, thumb is field). Since the loop is symmetric (centered at the origin of the \(x-y\) plane, the \(x-y\) components will cancel, and there will be a net field only in the \(z\) direction at any point along the \(z\) axis.

Then, finding \(d \vec{l}\), we may use cylindrical coordinates to find \(d \vec{l} = \hat{\alpha}_\phi a d \phi\). Then, \(\vec{R}\) can be expressed as \(\vec{R} = \hat{\alpha}_z \cdot z - \hat{\alpha}_r \cdot a\), and \(d \vec{l} \times \vec{R} = \hat{\alpha}_r a z d \phi + \hat{\alpha}_z a^2 d\phi\). 

Simplyfing, \(dH \to \hat{\alpha}_z\), and \(H = \int d \vec{H}\).

Solving, 
\[
	\vec{B} = \mu_0 \vec{H} = \hat{\alpha}_z \frac{mu_0 I \cdot a^2}{2(z^2 + a^2)^{\frac{3}{2}}} \ \mathrm{T}.
\]

When two current carrying conductors exert magnetic fields forces, there is an analogous relationship to that between two point charges; that is, the superposition principle is such that the forces contributed by each conductor sums with the force exerted by the other. In addition, there is a force on each conductor exerted by the other.

For two parallel conductors with current runnin gin the same direction along the \(z\) direction, we may find the force exerted on the second conductor, for example, as 
\[
  \vec{F}_2 = (\vec{I} \times \vec{B})L = -\hat{\alpha}_y \frac{\mu_0 I_1 I_2 l}{2\pi d}.
\]
Similarly, 
\[
  \vec{F_1} = \hat{\alpha}_y \frac{\mu_0 I_1 I_2}{2\pi d}.
\]
The takeaway here is there is an \textit{attractive} force between the wires. 

In terms of Maxwell's equations, magnetostatics can be reduced to 
\begin{align*}
	& \nabla \cdot \vec{D} = \rho_v\\
	& \nabla \times \vec{E} = 0 \\
	& \nabla \cdot \vec{B} = 0.
\end{align*}

Moving on, Ampere's Law relates the static magnetic field to the source of the field (current), and states that the line integral of the magnetic field \(\vec{H}\) around a closed path is equal to the current flowing through the surface bounded by that path. Mathematically, 
\[
	\oint_C \vec{H} \cdot d \vec{l} = I.
\]
In terms of magnetic flux density \(\vec{B}\), 
\[
  \oint_C \vec{B} \cdot d \vec{l} = \mu_0 I.
\]

Using Ampere's Law, we can solve the infinitely long conductor problem from before much more simply. 

Using the same conductor, infinitely long and lying along the \(z\) axis, we may find the magnetic field intensity at a distance \(r\) in the \(x-y\) plane. First, we can assume that the magnetic field is in the \(\phi\) direction, so \(\vec{H} = \hat{\alpha}_\phi H\). Then, \(d \vec{l} = \hat{\alpha}_\phi r d \phi\). Therefore, Ampere's Law in this case reduces to 
\[
  \int_{0}^{2\pi} \hat{\alpha}_\phi H \cdot \hat{\alpha}_\phi r d\phi = H\cdot r \cdot 2\pi = I.
\]

This agrees with our prior assessment, that 
\[
  H = \frac{I}{2\pi r}.
\]

As another example, a long straight wire of radius \(a\) carries a current \(I\) that is uniformly distributed over its cross section. Find the magnetic field \(\vec{H}\) a) inside the wire and b) outside the wire.

Again, we may center the wire at the origin of the \(x-y\) axis to solve more easily for \(\vec{H}\). Applying Ampere's Law again, if we create a closed loop inside the wire with radius \(r \leq a\), \(H 2\pi r = \frac{I}{\pi a^2} \cdot \pi r^2\), using the current density. Simplifying, 
\[
	\vec{H} = \hat{\alpha}_\phi\frac{I \cdot r}{2 \pi a^2} \ \mathrm{\frac{A}{m}}.
\]

Outside the loop, 
\[
	\vec{H} = \hat{\alpha}_\phi	\frac{I}{2 \pi r}.
\]

Continuing applications of Ampere's Law, find the magnetic flux density inside a closely wound toroidal coil with \(N\) turns of coil carrying a current \(I\). Also find \(\vec{B}\) for \(r<a\) and \(r>b\). 

Beginning with \(r<a\), we may use Ampere's Law again, 
\[
  \oint \vec{B} \cdot d \vec{l} = \mu_0 I.
\]
Using a closed loop, \(B \cdot 2\pi r = \mu_0 \cdot I\). Since \(I=0\); that is, there is no current flowing inside the toroid's inner radius, there is no magnetic field in the region \(r<a\), and \(\vec{B} = 0\).

Moving on to \(a<r<b\), we may once again assume that \(B\cdot 2\pi r = \mu_0 \cdot I\). What is the total current in the region \(a < r < b\)? Based on the number of turns, \(I_{\mathrm{total}} = I \cdot N\). Therefore, \(\vec{B} = \hat{\alpha}_\phi\ \frac{\mu_0 N I}{2\pi r}\). 

Finally, outside the toroid, in the region \(r > b\), \(I=0\) and \(\vec{B}=0\). 

\textbf{Note:} Double-check Ampere's Law and how it works, doesn't seem to be by \textit{enclosing} current, like Gauss's.

Revisiting Maxwell's equations, 
\begin{align*}
	& \nabla \cdot \vec{D} = \rho_v \\
	& \nabla \times \vec{E} = 0 \\
	& \nabla \cdot \vec{B} = 0 \\
	& \nabla \times \vec{H} = \vec{J}.
\end{align*}
The last equation correspond's to Ampere's Law, with a slightly different formulation. Integral of current density is current, duh. 

\textbf{Note:} Revisit Stokes' Theorem, and relationship to current/current density.

Static magnetic fields, like static electric fields, have boundary conditions. As a review of boundary conditions, we'd like to know how the magnetic field \(\vec{H}\) and the magnetic flux density \(\vec{B}\) change across a boundary with a difference in magnetic permeability.

The relevant equations are
\begin{align*}
	& B_{1n} = B_{2n} \\
	& \mu_1H_{1n} = \mu_2H_{2n} \\
	& \hat{n} \times (\vec{H}_1 - \vec{H}_2) = \vec{J}_s \\
	& H_{1t} - H_{2t} = J_s \\
\end{align*}

For most cases, \(J_s = 0\) and so \(H_{1t} = H_{2t}\).

\subsection{Inductors and Inductance}

Defining an \textbf{inductor}, an inductor is an energy storage device that stores energy in a magnetic field (similar to a capcitor, which stores energy in an electric field). As a reminder, \[C = \frac{Q}{V}.\]

For inductance, 
\[
  L = \frac{\Lambda}{I},
\]
that is, flux linkage in webers over current. The energy stored in an inductor may be described as 
\[
  W_m = \frac{1}{2}LI^2.
\]
What is flux linkage \(\Lambda\)? Flux linkage of an inductor is the total magnetic flux that links the current; similar to magnetic flux, but if you have \(N\) turns, like in the toroid example, the total magnetic flux \(\Phi\) that passes through those windings is the flux linkage.

Mathematically, 
\[
  \Phi = \int\int_S \vec{B} \cdot d \vec{s}.
\]
For \(N\) turns, 
\[
  \Lambda = N \cdot \Phi.
\]

Common types of inductors:

Solenoids have cylindrical shaped frames and have \(N\) turns of current-carrying wire wound around them. Deriving their inductance, 
\[
  L = \frac{\mu N^2 \pi a^2}{l},
\]
where \(a\) is the radius and \(l\) is the length. 

Toroids have donut-shaped frames with \(N\) turns of current-carrying wire wound around them. Deriving their inductance, 
\[
  L = \frac{\mu N^2 a^2}{2 \cdot r_0},
\]
where \(r_0\) is the average? radius and \(a\) is??

For example, if the magnetic flux \(\Phi\) produced by a given current links that same current, then the resulting inductance is self-inductance, but if it links the current in another circuit, then the resulting inductance is mutual inductance (requires two separate circuits/structures).

Self-inductance:
\[
	L = \frac{\Lambda}{I} = \frac{\mu}{I} \int \int_s \vec{B} \cdot d \vec{s}.
\]

Mutual inductance: 
\[
	L_{12} = \frac{\Lambda_{12}}{I_{1}} = \frac{N_2}{I_1} \int \int_{s_2} \vec{B}_1 \cdot d \vec{s}_{s_2}.
\]

The procedure for determining inductance is as follows:
\begin{enumerate}
	\item Choose a proper coordinate system
	\item Assume a current \(I\) in the conducting wire
	\item Find \(\vec{B}\) from \(I\). (i.e., use Ampere's Law)
	\item Find \(\Phi\) (magnetic flux) as
		\[
			\Phi = \int \int_S \vec{B} \cdot d \vec{s}
		\]
	\item Find flux linkage as
		\[
		  \Lambda = N \cdot \Phi
		\]
	\item Finally, 
		\[
		  L = \frac{\Lambda}{I}.
		\]
\end{enumerate}
As an example, develop an expression for the inductance per unit length of a coaxial line with inner and outer conductors of radii \(a\) and \(b\) and an insulating material of \(\mu\).

Using Ampere's Law, we can find the field that exists in the region \(a < r < b\) as 
\[
  \vec{B} = \hat{\alpha}_\phi \frac{\mu I}{2 \pi r}.
\]

Calculating the flux, 
\[
  \Phi = \int \int_s \vec{B} \cdot d \vec{s} = \int_0^l \int_a^b \hat{\alpha}_\phi \frac{\mu I}{2 \pi r} \cdot \hat{\alpha}_\phi \ dr \ dz.
\]
Therefore, 
\[
	\Phi = \frac{\mu I}{2 \pi} \cdot l \int_a^b \frac{1}{r} \ dr = \frac{u I l}{2 \pi} \ln{(\frac{b}{a})}.
\]
Finally, 
\[
  \Lambda = N \cdot \Phi = \Phi,
\]
and 
\[
	L = \frac{\Lambda}{I} = \frac{\mu l}{2 \pi} \ln{(\frac{b}{a})}.
\]
Therefore, self-inductance in this case is a function of permeability and inner/outer radius, primarily, if we divide by \(l\) to find inductance per unit length.

Mathematically, for a coaxial cable,
\[
	L' = \frac{\mu}{2 \pi} \ln{(\frac{b}{a})}.
\]

As a final example, find the mututal inductance between a conducting rectangular loop and very long, straight wire.

Assuming a current of \(I\) on the straight wire, we will first find \(\vec{B}\) using Ampere's Law in the form \(\vec{B} = \hat{\alpha}_\phi \frac{\mu I}{2 \pi r}\). Then, derive \(\Phi\) (\textbf{Note:} Review this in the book, brain cells fried).

More examples of inductance.

For example, \(N\) turns of wire are tightly wound on a toroidal frame of a rectangular cross section with dimensions shown below. Assuming the permeability of the medium to be \(\mu_0\), find the self inductance.

Recall that inductance is defined as 
\[
  L = \frac{\Lambda}{I},
\]
\[
	\Lambda = N \cdot \Phi,
\]
and 
\[
  \Phi = \int \int_S \vec{B} \cdot d \vec{s}.
\]

Assuming a current of \(I\) in the wire, we can find \(\vec{B}\) from \(I\); that is, the magnetic flux density. Using Ampere's Law,
\[
	\oint_c \vec{B} \cdot d \vec{l} = \mu_0 I_{\mathrm{tot}}.
\]
Multiplying current \(I\) by number of turns \(N\), 
\[
  \vec{B} = \hat{\alpha_\phi} \ \frac{\mu_0 N I}{2 \pi r}.
\]
Now knowing the magnetic flux density \(\vec{B}\), we may find the magnetic flux \(\Phi\). Using the aforementioned formula, 
\[
	\Phi = \int_{0}^{h} \int_{a}^{b} (\hat{\alpha_\phi} \ \frac{\mu_0 N I}{2\pi r}) \cdot (\hat{\alpha_\phi \ dr \ dz}) = \frac{\mu_0 N I h}{2\pi} \ \ln({\frac{b}{a}}).
\]
Now, with \(\Phi\), we can find flux linkage \(\Lambda\), using the aforementioned formula,
\[
	\Lambda = N \cdot \Phi = \frac{\mu_0 N^2 I h}{2\pi} \ln{(\frac{b}{a})}.
\]
Finally, inductance \(L\) can be found as
\[
	L = \frac{\Lambda}{I} = \frac{\mu_0 N^2 h}{2\pi} \ \ln{(\frac{b}{a})} \ \mathrm{H}.
\]
\end{document}

